(1) 纯转动能为 \(K=I\omega^2/2\)，恒输出功率条件 \(P=dK/dt=I\omega\dot\omega\)。积分得 \(\omega^2(t)=\omega_i^2+2Pt/I\)，因此 \(\omega(t)=\sqrt{\omega_i^2+2Pt/I}\)，而非恒角加速度线性增长。
(2) 电机输出转矩 \(\tau=P/\omega(t)\)，由 \(I\dot\omega=\tau\) 得 \(\dot\omega=P/(I\omega(t))\)。当形式上令 \(\omega_i=0\) 时 \(t=0\) 的转矩趋无穷，真实电机必须先以有限转矩起步，之后才可能进入近似恒功率区域；此模型需明确非零初始转速。
(3) 薄环外缘速度为 \(v=R\omega\)，所以 \(\omega_{\max}=v_{\max}/R\)。可增加的转动能 \(\Delta K_{\max}=\tfrac12MR^2[(v_{\max}/R)^2-\omega_i^2]=\tfrac12M(v_{\max}^2-R^2\omega_i^2)\)，加速时间 \(t_{\max}=\Delta K_{\max}/P\)，要求 \(R\omega_i\le v_{\max}\)。
(4) 恒功率再生制动给 \(dK/dt=-P_b\)，因此 \(\omega^2(t)=\omega_{\rm start}^2-2P_bt/I\)。达到目标角速 \(\omega_f\) 的时间为 \(t_b=I(\omega_{\rm start}^2-\omega_f^2)/(2P_b)\)，恰好满足 \(P_bt_b=K_{\rm start}-K_f\)；若制动到零速，恒功率阶段末端同样需要以有限转矩控制替代。
